\setproblem{1999}

\begin{frame}[label=p1]{최대최소 (\prob)}
\small
\begin{block}{문제}
$N \times N\ (1 \le N \le 250)$ 행렬이 주어진다. 각 원소는 $250$ 이하의 음이 아닌 정수이다. \\

크기 $B \times B\ (1 \le B \le N)$의 부분행렬에 대해,
\textbf{최댓값과 최솟값의 차이}를 구하는 질의가
$K\ (1 \le K \le 100\,000)$개 주어진다. \\

모든 질의는 같은 $B$를 사용하며, 각 질의에서는 부분행렬의 왼쪽 위 좌표 $(i,j)$가 주어진다.
\end{block}
\end{frame}


\begin{frame}{예제 입력 / 출력}
\huge
\[
\begin{array}{|c|c|c|c|c|}
\hline
5 & 1 & 2 & 6 & 3 \\ \hline
1 & 3 & 5 & 2 & 7 \\ \hline
7 & 2 & 4 & 6 & 1 \\ \hline
9 & 9 & 8 & 6 & 5 \\ \hline
0 & 6 & 9 & 3 & 9 \\ \hline
\end{array}
\]
\end{frame}

\begin{frame}{$\mathcal{O}(N^4)$ 풀이}
가능한 모든 $B \times B$ 크기의 부분 행렬의 최솟값, 최댓값을 구해놓습니다.
\[
\begin{array}{|c|c|c|c|c|}
\hline
5 & 1 & 2 & 6 & 3 \\ \hline
1 & 3 & 5 & 2 & 7 \\ \hline
7 & 2 & 4 & 6 & 1 \\ \hline
9 & 9 & 8 & 6 & 5 \\ \hline
0 & 6 & 9 & 3 & 9 \\ \hline
\end{array}
\]
\end{frame}

\begin{frame}{$\mathcal{O}(N^3)$ 풀이}
\small
$B \times B$ 크기의 부분 행렬의 최대(소)값은 행, 열 단위의 최솟값을 이용하여 구할 수 있습니다.
$p_{i,j} = \max_{0 \le k < B} a_{i, j+k}$가 되도록 각 행마다 $\mathcal{O}(N \times B)$ 시간에 만들 수 있습니다.

\[
\begin{array}{|c|c|c|c|c|}
\hline
5 & 1 & 2 & 6 & 3 \\ \hline
1 & 3 & 5 & 2 & 7 \\ \hline
7 & 2 & 4 & 6 & 1 \\ \hline
9 & 9 & 8 & 6 & 5 \\ \hline
0 & 6 & 9 & 3 & 9 \\ \hline
\end{array}
\qquad
\begin{array}{|c|c|c|c|c|}
\hline
5 & 6 & 6 & - & - \\ \hline
5 & 5 & 7 & - & - \\ \hline
7 & 6 & 6 & - & - \\ \hline
9 & 9 & 8 & - & - \\ \hline
9 & 9 & 9 & - & - \\ \hline
\end{array}
\]

이제 이 방법을 열마다 적용할 수 있습니다. $q_{i, j} = \max_{0 \le k < B} p_{i+k, j}$가 되도록 각 열마다 $\mathcal{O}(N \times B)$ 시간에 만들 수 있습니다.
\end{frame}

\begin{frame}{$\mathcal{O}(N^2)$ 풀이}
\small
위에서 구간 최대 최소를 구하는 방법을 자료구조로 최적화 할 수 있습니다. 시간 복잡도는 $\mathcal{O}(N\log{N})$입니다.
\\

그런데, 고정된 $B$에 대해서는 $\mathcal{O}(N)$ 시간에 구하는 방법이 있습니다. 최소값을 덱에 증가하는 순서로 저장하는 방식입니다.

\begin{enumerate}
\item $a_1, a_2, \ldots, a_N$를 순서대로 확인하면서 덱의 뒤에 넣거나 덱의 앞 원소를 제거합니다.
\item 덱에는 최소값이 될 수 있는 후보를 넣습니다. 단, 더 작은 원소가 앞서도록 유지해야합니다.(단조성)
\item $a_i, a_{i + 1}, \cdots, a_{i + B - 1}$를 확인하고 덱에 넣었다면 덱의 맨 앞 원소가 $[i, i + B - 1]$ 구간의 최소값이 됩니다.
\item 매번 $a_i$를 덱에 넣을 때마다 덱의 맨 뒤 원소가 더 크다면 제거해고 나서 $a_i$를 넣습니다.
\item 매번 덱 앞 원소가 $[i, i + B - 1]$ 구간에서 벗어난다면 제거해야 합니다.
\end{enumerate}

\end{frame}

\begin{frame}[fragile, allowframebreaks]{원찬혁 소스 코드}
\inputminted{java}{../src/\prob/chanhyeok/Main.java}
\end{frame}

\begin{frame}[fragile, allowframebreaks]{원찬혁 소스 코드 2}
\inputminted{java}{../src/\prob/chanhyeok/Main2.java}
\end{frame}

\begin{frame}[fragile, allowframebreaks]{손현준 소스 코드}
\inputminted{java}{../src/\prob/hyunjoon/P1999.java}
\end{frame}

\begin{frame}[fragile, allowframebreaks]{김준형 소스 코드}
\inputminted{java}{../src/\prob/joonhyeong/Main.java}
\end{frame}

\begin{frame}[fragile, allowframebreaks]{서민종 AccumulationTwo 소스 코드}
\inputminted{java}{../src/\prob/minjong/AccumulationTwo.java}
\end{frame}

\begin{frame}[fragile, allowframebreaks]{서민종 TreeMapSolution 소스 코드}
\inputminted{java}{../src/\prob/minjong/TreeMapSolution.java}
\end{frame}

\begin{frame}[fragile, allowframebreaks]{서민종 VeryLong 소스 코드}
\inputminted{java}{../src/\prob/minjong/VeryLong.java}
\end{frame}

\begin{frame}[fragile, allowframebreaks]{정의찬 소스 코드}
\inputminted{java}{../src/\prob/uichan/1999.java}
\end{frame}
